\documentclass{article}
\usepackage{amsmath, amssymb}

\begin{document}

\section*{Q1) Range Sum Queries}

\textbf{Problem:} \\
You are given an array of $N$ integers. You are asked $Q$ queries. Each query provides two indices $L$ and $R$. For each query, efficiently compute:
\[
\sum_{i=L}^{R} A[i]
\]

\textbf{Idea:} \\
Use prefix sums:
\[
P[i] = A[1] + A[2] + \dots + A[i]
\]
Then each query can be answered in $O(1)$:
\[
\text{Sum}(L,R) = P[R] - P[L-1]
\]

\section*{Q2) Longest Subarray with Sum $\leq X$}

\textbf{Problem:} \\
Given an array of positive integers and a target value $X$, find the length of the longest contiguous subarray whose sum does not exceed $X$.

\textbf{Idea:} \\
Use two pointers (sliding window). Expand right pointer and shrink left pointer when sum exceeds $X$. Maintain maximum valid window length.

\section*{Q3) Factory Machines}

\textbf{Problem:} \\
You have $N$ machines. The $i$-th machine takes $t_i$ time to produce one product. All machines work in parallel. Find the minimum time required to produce exactly $T$ products.

\textbf{Idea:} \\
Use binary search on time $T_{\text{time}}$. For a fixed time $t$, total products produced is:
\[
\sum_{i=1}^{N} \left\lfloor \frac{t}{t_i} \right\rfloor
\]
Check feasibility and minimize using binary search.

\section*{Q4) Bit Manipulation: Sum of Pairwise XOR}

\textbf{Theory:} \\
Numbers are represented in binary. Bitwise operations:
\begin{itemize}
    \item AND ($\&$): 1 if both bits are 1
    \item OR ($|$): 1 if at least one bit is 1
    \item XOR ($\oplus$): 1 if bits differ
    \item Left shift ($<<$): multiply by $2^k$
    \item Right shift ($>>$): divide by $2^k$
\end{itemize}

\textbf{Problem:} \\
Given an array of $N$ integers, compute:
\[
\sum_{1 \le i < j \le N} (A[i] \oplus A[j])
\]

\textbf{Idea:} \\
For each bit position, count number of $1$s and $0$s. Contribution:
\[
\text{bit contribution} = (\text{ones}) \cdot (\text{zeros}) \cdot 2^b
\]
Sum over all bit positions gives $O(N \log A)$ solution.

\section*{Q5) Greedy: Interval Scheduling}

\textbf{Theory:} \\
A greedy algorithm makes the locally optimal choice at each step to reach a global optimum.

\textbf{Problem:} \\
Given $N$ intervals $(s_i, e_i)$, select the maximum number of non-overlapping intervals.

\textbf{Idea:} \\
Sort intervals by end time. Always pick the interval with the earliest finishing time that is compatible.

\section*{Advanced Topic 1: Ternary Search}

\textbf{Theory:} \\
Ternary search finds the maximum or minimum of a unimodal function (increasing then decreasing).

\textbf{Idea:} \\
Split range $[L,R]$ into three parts using midpoints $m_1$ and $m_2$. Compare $f(m_1)$ and $f(m_2)$ and discard one-third each step.

\section*{Advanced Topic 2: Maximum Subarray Median}

\textbf{Theory:} \\
We use binary search on the answer and prefix sums.

Transform array using a guess $M$:
\[
A[i] =
\begin{cases}
+1 & \text{if } A[i] \ge M \\
-1 & \text{otherwise}
\end{cases}
\]

\textbf{Problem:} \\
Find the maximum median of any subarray with length at least $K$.

\textbf{Idea:} \\
Binary search on median $M$. Use prefix sums to check if a valid subarray with positive transformed sum exists.

\end{document}